\section{Absolute phase and the regular sewing density}\label{sec:APS}

Let all spine fibres of the boundary local system \(E\) act trivially, so the same is true of all spine sections. The system can be the rank-three adjoint of a flat \(SU(2)\) connection with central spine fibres. Write \(I_v\) for the invariant subspace of the entire local system at \(v\), and \(R_e\) for its orthogonal or unitary edge transport. Put \(Q_I\) on \(\bigoplus_v I_v\) with diagonal \(D_v I\) and edge pairing \(-\epsilon_e\langle u_v,R_eu_w\rangle\). For complex rank-one systems use \(\rho_E=\eta({\cal D}_E)-\operatorname{rank}(E)\eta({\cal D}_{\rm triv})\).

\begin{proposition}[Deleted-trace formula]\label{prop:aps}
For the oriented central local system above,
\begin{equation}\label{eq:deletion}
 \rho_E(M) = \operatorname{rank}(E)\sig Q-\sig Q_I+
             \sum_{\text{caps}}\rho_{E|_{\partial N}}(\partial N),
\end{equation}
where \(N\) is the actual corresponding terminal disk bundle with its trace orientation. This is an identity of real numbers also with nullities.
For the opposite order-two pair on every vertex, on the branch with all fibres \(-1\), let \(u_v\) be the unit cone axis and \(B_{vw}=\epsilon_e\langle u_v,R_eu_w\rangle\) across an edge. Then
\begin{equation}\label{eq:pairedaps}
 \CS=l/2\pmod1,\qquad \rho_{\rm ad}=\sig B-3\sig A_\epsilon,\qquad l=|E(G)|.
\end{equation}
\end{proposition}

\begin{proof}
Use the actual disk-bundle plumbing trace, with terminal weights \(p\) and edge intersections \(-\epsilon_e\) in the stated peripheral convention. The latter sign follows by exchanging the fibre and base disk coordinates: the section bounding the deleted base disk has the opposite orientation from its base complex angle. Delete the interiors of the small terminal disk bundles \(N\). Locally the surviving central bundle retracts to the punctured sphere; the remaining terminal region retracts to \(\partial N\) and overlaps the central region in a meridian circle up to homotopy. Ordinary graph-edge overlaps are contractible and no two such overlaps meet. As the boundary fibre is trivial in \(E\), the cone representation extends across the local trace by the order-\(|p|\) lens relation. Graph stable transports extend as well, since the edge overlaps are contractible. In the negative-central order-two case the fundamental representation need not extend; the adjoint does extend, with each order-two cone generator a half-turn.

The terminal lens pieces have no degree-one or degree-two homology with these coefficients, by finite transfer over characteristic zero. At a vertex the local Mayer--Vietoris group in degree two is therefore the kernel of the peripheral-circle maps into the punctured base, equivalently \(H_2(P_v,\partial P_v;E)\cong I_v\) by twisted duality (use \(H_2(S^2;E)\) directly when no cap is removed). For an invariant coefficient the representative is the central sphere with its rational terminal corrections \(S_v-\sum_jS_{vj}/p_{vj}\). Its square is \(D_v\); at an edge it uses the unchanged central chart. Intersecting these representatives pairs their transported coefficients with sign \(-\epsilon_e\). Thus the ordinary and twisted middle signatures of the deleted trace are exactly \(\sig Q,\sig Q_I\), with radicals omitted, including when the graph has cycles. APS subtraction in the positive mixed-operator convention gives \(\rho_E(\partial W^\circ)=\operatorname{rank}(E)\operatorname{Sign} W^\circ-\operatorname{Sign}_E W^\circ\)~\cite{APS,KKR}. The extra boundaries are \(-\partial N\), proving \eqref{eq:deletion}. For opposite order-two caps \(I_v=\R u_v\), hence \(Q_I=-B\); the two lens systems are isomorphic with opposite orientations. This proves the rho equality.

For the action on this branch, deform the stable transports to a common axial frame; the cap and central seam relations remain flat. There is an abelian endpoint with spine eigenline coordinates \(-1/2\), plus-leaf coordinates \(1/4\), and minus-leaf coordinates \(s_v/2-1/4\), where \(s_v=\sum_{w\sim v}\epsilon_{vw}\). Let this vector be \(\alpha\), and \(L\) the integral trace matrix (cap incidence can be made \(+1\)). Then \(L\alpha\) is integral: its entries are \(s_v\) on the spine, zero on the plus leaves, and \(-s_v\) on the minus leaves. Thus the eigenline class has the rational relative lift \(\alpha\); its square is \(\alpha^{\mathsf T}L\alpha\) modulo one, valid even if the matrix is singular because the lift exists. The free graph phases can be extended on the trace and do not affect the square. The spine contribution \(-\sum_v s_v/2\) is integral. The remaining term, up to overall sign, is \(\sum_v s_v(2s_v-1)/4\equiv l/2\pmod1\), since \(s_v(s_v-1)\) is even and \(\sum s_v=2\sum_e\epsilon_e\). Eigenline transgression as in Section~\ref{sec:H} and flat continuity give the stated CS.
\end{proof}

\begin{theorem}[Actual one-loop comparison for closed odd cycles]\label{thm:cycle}
For opposite order-two caps and a closed odd spine cycle of length \(n\), let \(\sigma=\sig A_\epsilon\). The all-negative-fibre locus is smooth at its irreducible points with \(\det B\ne0\). There the adjoint Reidemeister half-density is \(c_G|\det B|^{-1/2}d\mu\), with \(c_G>0\) constant and \(d\mu\) quotient Haar. The full unmarked quantum asymptotic, normalized by \(r^n\), is
\begin{equation}\label{eq:cycleaps}
 e^{2\pi\ii k\CS}\frac{2^{n/2}}{\pi^n}\,2^{-n/2}
   \E_\mu\left[\frac{e^{-\pi\ii\rho_{\rm ad}/4}}{\sqrt{|\det B|}}\right]+o(1).
\end{equation}
The expectation is absolutely convergent. In particular it gives \(Z_r/r^n=(-1)^{rn}e^{3\pi\ii\sigma/4}2^{n-1}/\pi^n+o(1)\).
\end{theorem}

\begin{proof}
At each vertex the two cone quaternions are \(u_v\) and a specified \(\pm u_v\); all marked torus holonomies are \(-1\). After fixing the axes, the edge conjugators are free, with one axial \(U(1)\) gauge action per vertex. At a regular point there are no additional nearby fibre defects. Indeed the centrality and two cap square relations require the axial projection of the logarithm of the sign-stripped product of sections to vanish. Its derivative in the fibre angle defects is \(B\); at zero defect these equations hold for all nearby axis and transport choices. The implicit function theorem and the existing flat all-central family therefore identify the local locus before the axial quotient. At irreducible points the axial action has only the finite central stabilizer (trivial in the effective action), so the quotient is smooth there. No such smoothness is claimed at reducibles: aligned transports may have \(\det B\ne0\) and still have an axial stabilizer. Reducible transports have \(T_e=\pm1\) on each edge, since all local cone axes must lie on one transported invariant line, and are consequently Haar-null.

For the torsion density, right-trivialized cocycles give \(H^0\) the axis line at each block. Their boundary \(H^1\) traces are a common fibre variation \(p_vu_v\) and section variations \(s_{ve}\) subject to \(\sum_e\langle u_v,s_{ve}\rangle=0\): the two square relations determine longitudinal cap variations, and their product determines the second transverse variation after killing the first by a coboundary. This restriction is injective and has rank \(3d_v\). \(H^2\) injects into boundary area-coordinates with the same summed axial condition by the pair sequence and duality. Put \(P\) equal to the space of all incidence section vectors and \(Ua=(a_vu_v)_{ve}\). The degree-one gluing map has the form \((s,p)\mapsto s-Kp\) on \(U^{\mathsf T}s=0\), where \(Kp\) is the swapped transported fibre variation. Augmentation by \(U^{\mathsf T}s\) on all of \(P\) leaves determinant \(|\det U^{\mathsf T}K|=|\det B|\). These volumes are fixed block quotient volumes. In degree zero the seam map is \(J={\cal A}U\), where \({\cal A}s\) is the difference of transported incidence vectors. It is injective at an irreducible. In degree two the seam map is the corresponding sum \({\cal C}\) on \(\ker U^{\mathsf T}\); \(\ker{\cal C}=\operatorname{im}{\cal A}^{\mathsf T}\) and \(U^{\mathsf T}{\cal A}^{\mathsf T}=J^{\mathsf T}\). For tangent lifts \(y\) to the seam space, put \(c=|\det[J,y]|\). The Poincaré-dual degree-two basis is \({\cal A}^{\mathsf T}b\), where \(b\in\ker J^{\mathsf T}\) is dual to \(y\); the duality follows by integrating the \(H^1\) connecting representative along the seam collar. Appending to it a corrected one-sided lift of the degree-two seam map gives determinant \(c^{-1}\) by the dual-basis identity. Thus multiplicativity of cochain torsion with dual bases gives \(c^2/|\det B|\), up to constant block determinants (an isolated isomorphism from degree one to degree two contributes the inverse determinant). See~\cite[Prop.\ 3.3]{Massuyeau}. The quotient Haar volume in tangent coordinates is proportional to \(c\), by the axial gauge columns \(J\). Taking a square root proves the density assertion; no absolute normalization of a functional measure is implied.

Independently, let \(d=r-a\). The exact cap identity of Section~\ref{sec:H} and the edge reflection sign on the supported odd defects give
\[
 Z_r=(-1)^{rn}\kappa_r^\sigma\frac{2^{3n/2}}{\pi^n}r^{n/2}
 \sum_{\boldsymbol d\ {\rm odd}<r}\,
  \prod_v \left[\cos^2(\pi d_v/(2r))
             \left(\frac{\pi d_v/r}{\sin(\pi d_v/r)}\right)^2\right]
  \prod_{vw}\sinc(\pi d_vd_w/r).
\]
The bracket is bounded on the entire interval including the upper end. At \(d_v=\sqrt r\,t_v\), each odd lattice gives density \(\sqrt r/2\), yielding coefficient \(2^{n/2}\pi^{-n}{\cal J}\) with \({\cal J}=\int_{t>0}\prod_{vw}\sinc(\pi t_vt_w)\,d^nt\). Uniform tails follow on dyadic boxes \(t_v\asymp2^{j_v}\): the volume times sinc bound is at most \(C2^{-\sum_{vw}|j_v+j_w|/2}\), since the graph has degree two. The edge-sum linear map is invertible on an odd cycle, so these bounds sum at all scales. This is a whole closed limit, not an inserted corner.

The normalized Haar marks \(T_e=\langle u_v,R_eu_w\rangle\) are independent uniform on \([-1,1]\), also around graph cycles: the unconstrained edge conjugators are product Haar before dividing by the constant gauge volume. Write each sinc as this orbital integral; with positive Gaussian damping Fubini holds and the sign-symmetric edge law extends the orthant with factor \(2^{-n}\). Its Gaussian gives \({\cal J}=2^{-n/2}\E_\mu e^{-\pi\ii\sig B/4}/\sqrt{|\det B|}\), choosing the minus sign by the symmetry \(B\mapsto-B\) before the limit. The determinant envelope is integrable: \(\det B=2\prod_e\epsilon_eT_e\), only the two full directed cycles contribute, and \(\E|T_e|^{-1/2}=2\). The signatures are the opposite values \(+1,-1\) on the two sign classes (deform magnitudes to one and diagonalize the cycle), so \({\cal J}=2^{n/2-1}\). Finally combine \(\kappa_r^\sigma\) with \eqref{eq:pairedaps} and the independent action \(e^{2\pi\ii k\CS}=(-1)^{rn}\). This proves \eqref{eq:cycleaps}. The rho identity holds beyond this odd-cycle example but does not itself show integrability or a nonzero leading coefficient on other graphs.
\end{proof}
